%%%%%%%%%%  scan idx 79  |  folio 66  %%%%%%%%%%
\SGAPrintPage{66}%
Formula (iv) expresses that $i_!$ is a ring homomorphism when we introduce on $\underline K(Y)$ a new multiplicative structure, by the definition $x_N x'=xx'\lambda_{-1}(N)$ (cf.\ Chapter I, \S 4). It is completed by the formula:
\begin{enumerate}[label=(\roman*)]
\setcounter{enumi}{4}
\item $i_!(\lambda^p(N,x))=\lambda^p(i_!(x))$ for $p\geq1$,
\end{enumerate}
which means that $i_!$ defines a $\lambda$-homomorphism from $\underline K(Y)_N$ to $\underline K(X)$. This formula (v), certainly true in all cases, has for the moment been proved only in characteristic $0$, since it follows from the conjunction of theorems 1.4 and 2.8.

We have analogous formulas for $i_*$ and $i^*$:
\begin{enumerate}[label=(\roman* bis)]
\item $i^*$ is a homomorphism of graded rings.
\item $i_*(x.i^*(y))=i_*(x).y$, whence in particular $i_*(i^*(y))=\xi' y$, with $\xi'=i_*(1)$,
\item $i^*(i_*(x))=x.(-1)^q N^q$, whence in particular $i^*(\xi')=(-1)^q N^q$,
\item $i_*(x)i_*(x')=i_*(xx'(-1)^q N^q)$,
\item $i_*$ increases degrees by $q$ units.
\end{enumerate}

In these formulas, $q$ is the codimension of $Y$ in $X$, and we put $N^i=C^i(N)$. Only formulas (iii bis) and (iv bis) require a proof; I do not know whether they are well known (or even whether they are true!!!), and merely note that they follow by reduction of formulas (iii) and (iv), \emph{provided that we work in} $G(\underline K(Y))$ \emph{and} $G(\underline K(X))$ \emph{and not in} $A(Y)$ \emph{and} $A(X)$. We then use the fact that $i_*$ and $i^*$ arise by passage to the quotient from $i_!$ and $i^!$, noting that $i_!$ shifts the filtration by $q$ units.

%%%%%%%%%%  scan idx 80  |  folio 67  %%%%%%%%%%
\SGAPrintPage{67}%
N.B. All the ``universal'' formulas relating $i_!$ and $i^!$ (resp.\ $i_*$ and $i^*$) are formal consequences of the preceding ones, as one can convince oneself without much difficulty.

From formula (v) follows the formula
\[
 i_!(\gamma^p(N,x))=\gamma^p(i_!(x)).
\]
The conjunction of theorem 2.14 and proposition 1.5 shows that $\gamma^p(N,x)$ has, for $p=q+j$, filtration $\geq j$; denoting by $G_p(N,x)$ its class in $G^j(\underline K(Y))$, theorem 2.16 shows that we have
\[
 i_*(G_p(N,x))=c^p(i_!(x)).
\]
Finally, proposition 1.5 together with theorem 2.14 allows us to express $G_p(N,x)$ as a polynomial with integer coefficients in the $c^i(N)$, the $c^i(x)$ and $\varepsilon(x)$; we deduce the theorem:

\sgaStableTarget{sga6.e00.thm.2-17}\textbf{Theorem 2.17.} \emph{Assume that the field $k$ has characteristic $0$ and let $i:Y\longrightarrow X$ be an injection morphism of nonsingular quasi-projective varieties. We then have formulas (i) to (v), and, moreover, for every integer $j\geq0$, the formula:}
\sgaStableTarget{sga6.e00.eq.2-37}\begin{equation}\tag{2.37}
 c^{q+j}(i_!(x))=i_*\left(\frac{\widetilde c(x)*\lambda_{-1}(\widetilde c(N))}{(-1)^qN^q}\right)^{(j)}.
\end{equation}

The right-hand side of formula (2.37) is an abbreviated way of denoting the polynomial in the components of $\widetilde c(x)$ and $\widetilde c(N)$ described in proposition 1.5.

\sgaStableTarget{sga6.e00.cor.unnamed.occ5}\textbf{Corollary:} \emph{Under the preceding conditions, we have:}
\sgaStableTarget{sga6.e00.eq.2-38}\begin{equation}\tag{2.38}
 \ch(\widetilde c(i_!(x)))=i_*\left(\ch(\widetilde c(x))\,\mathcal T(\check N)^{-1}\right)
\end{equation}
\emph{where $\ch$ is the Chern homomorphism defined by formula} (1.26).

%%%%%%%%%%  scan idx 81  |  folio 68  %%%%%%%%%%
\SGAPrintPage{68}%
It suffices to use (2.37) written in the form:
\sgaStableTarget{sga6.e00.eq.2-37bis}\begin{equation}\tag{2.37 bis}
 \widetilde c(i_!(x))=1+i_*\left(\frac{\widetilde c(x)*\lambda_{-1}(\widetilde c(N))-1}{(-1)^qN^q}\right),
\end{equation}
the explicit formula (1.29), formula (iv bis), and finally (1.30).

We can write (2.38) in the form
\sgaStableTarget{sga6.e00.eq.2-38bis}\begin{equation}\tag{2.38 bis}
 \ch_X(F)=i_*\left(\ch_Y(F)\,\mathcal T(T_{X/Y})^{-1}\right),
\end{equation}
where $F$ is a coherent algebraic sheaf on $Y$ (considered on the left-hand side as a sheaf on $X$ zero outside $Y$), $\ch_X$ and $\ch_Y$ are the Chern homomorphisms taken respectively on $X$ and $Y$, and $T_{X/Y}$ is the normal bundle of $Y$ in $X$.

\sgaStableTarget{sga6.e00.rem.unnamed.occ5}\textbf{Remarks:} Unlike (2.38), formula (2.37) is a ``formula without denominators''. It has been proved (even in characteristic $0$) only by working in $G(\underline K(X))$ and not in $A(X)$, which may not be the same thing. It is certainly true in every characteristic and in $A(X)$. Let us note that we immediately verify that the images of the two sides under $i^*$ are the same, whence it follows, on applying $i_*$, that the product of their difference by $\xi'$ is zero. This suggests a completely different method of proof, by trying to obtain $Y$ as the transverse inverse image of a subvariety $Y'$ in a variety $X'$ such that, in dimension $\leq n=\dim X$, the class of $Y'$ in $A(X')$ is not a divisor of $0$\dots. Unfortunately, we cannot hope that $X'$ is a Grassmannian or something close to one.

We shall note that, if we disregard torsion elements, formula (2.38) is equivalent to (2.37).




% Additional macro definitions
\providecommand{\cS}{\mathcal S}
\providecommand{\cD}{\mathcal D}
\providecommand{\cE}{\mathscr E}
\providecommand{\cG}{\mathscr G}
\providecommand{\cL}{\mathscr L}
\providecommand{\cM}{\mathscr M}
\providecommand{\cP}{\mathscr P}
\providecommand{\cK}{\mathcal K}
\providecommand{\Db}{\mathrm D}
\providecommand{\Kb}{\mathrm K}
\providecommand{\Trc}{\operatorname{Tr}}
\providecommand{\ob}{\operatorname{ob}}
\providecommand{\Coh}{\operatorname{Coh}}
\providecommand{\Cone}{\operatorname{Cone}}
\providecommand{\im}{\operatorname{Im}}
\providecommand{\Id}{\operatorname{Id}}
\providecommand{\RR}{\mathrm{R.R.}}
\providecommand{\RRR}{\mathrm{R.R.R.}}
\providecommand{\GrR}{\mathrm{G}}
\providecommand{\pr}{\operatorname{pr}}


%%%%%%%%%%  scan idx 82  |  folio 69  %%%%%%%%%%
\SGAPrintPage{69}%
\sgaStableTarget{sga6.e00.sec.6.occ2}\section*{\S 6. The Riemann--Roch theorem}

Its statement is as follows: \emph{let \(f\) be a} proper \emph{morphism from a nonsingular quasi-projective variety \(Y\) to a variety \(X\) of the same kind. Then we have:}
\sgaStableTarget{sga6.e00.eq.2-39}\begin{equation}\tag{2.39}
 f_*\bigl(\ch_Y(x)\,\mathcal T(Y)\bigr)
 =
 \ch_X\bigl(f_!(x)\bigr)\,\mathcal T(X).
\end{equation}
(Both sides are in \(A(X)\otimes\mathbb Q\) and \(\mathcal T(X)=\mathcal T(C(T_X))\).)

It allows us to calculate, up to torsion elements, the Chern classes of \(f_!(x)\) from the Chern classes of \(x\), those of \(T_X\) and \(T_Y\), and the homomorphism \(f_*\). Or, again, it allows us to determine the homomorphism
\(f_!\) from \(\underline K(Y)\otimes\mathbb Q\) to \(\underline K(X)\otimes\mathbb Q\), when we identify these groups by means of the homomorphisms \(\ch_Y\) and \(\ch_X\) respectively with \(A(Y)\otimes\mathbb Q\) and \(A(X)\otimes\mathbb Q\) (if formula (2.35) is true!), using \(f_*\) and the Chern classes of \(T_X\) and \(T_Y\).

I know how to prove formula (2.39) only in characteristic \(0\)\SGAAuthorFootnote{*}{See below for a proof valid in every characteristic, based on a different principle (p.\,51).}, and as an equality in \(G(\underline K(X))\otimes\mathbb Q\). For this, since \(Y\) is locally closed in a projective space \(P\), we decompose \(f\) into the product of an injection morphism from \(Y\) to \(P\times X\), and a projection morphism from \(P\times X\) onto \(X\), and prove (2.39) for each of these two morphisms.

For an injection morphism, we see immediately that formulas (2.38) and (2.39) are equivalent (the characteristic \(0\) hypothesis entered only through theorem 2.8, which allows us to express the \(\lambda^i(\gamma(F))\), for a coherent algebraic sheaf \(F\), by \emph{local} cohomological invariants of \(F\)).

For a projection morphism \(f:P\times X\longrightarrow X\), we can, by proposition 2.13, assume that \(x\) is of the form \(y\otimes z\), with \(y\in\underline K(P)\) and
%%%%%%%%%%  scan idx 83  |  folio 70  %%%%%%%%%%
\SGAPrintPage{70}%
\(z\in\underline K(X)\).
Then the K\"unneth formula gives \(f_!(x)=\chi(y)z\), where, for a coherent algebraic sheaf \(F\) on \(P\), we put:
\[
  \chi(F)=\sum_{i\geq 0}(-1)^i\dim H^i(P,F).
\]
We are then immediately reduced to proving the formula
\sgaStableTarget{sga6.e00.eq.2-40}\begin{equation}\tag{2.40}
  \chi(F)=\varkappa_r\bigl(\ch_P(F)\,\mathcal T(P)\bigr)
\end{equation}
where \(\varkappa_r\) denotes the component of degree \(r=\dim P\). By proposition 2.13, we can assume that \(F=\underline{O}_Q\), where \(Q\) is a linear subvariety of \(P\). Then \(\chi(F)=1\), since it is the arithmetic genus of \(Q\); it remains to see that the right-hand side of formula (2.40) equals \(1\).

Indeed, if \(V\) is a hyperplane of \(P\) and we put \(\xi=\gamma_P(\underline{O}_V)\) and \(L=\ell(V)\), we shall have:
\[
  1-\xi=\gamma_P(E(-L))
\]
(where \(E(-L)\) is the vector bundle defined by \(-L\)), whence:
\[
  C_P(1-\xi)=[1,1-L],
\]
and consequently
\[
  \ch_P(1-\xi)=e^{-L},\qquad\text{whence}\qquad \ch_P(\xi)=1-e^{-L},
\]
and therefore
\[
  \ch_P(\xi^p)=(1-e^{-L})^p.
\]
Now, if \(p\) is the codimension of \(Q\) in \(P\), we have \(\xi^p=\gamma_P(Q)\); we deduce
\[
  \ch_P(F)\,\mathcal T(P)
  =
  \bigl(L/(1-e^{-L})\bigr)^{r+1}(1-e^{-L})^p
\]
and we must verify that the coefficient of \(L^r\) in this formal series in \(L\) equals \(1\), which follows from the characteristic property of the formal series \(L/(1-e^{-L})\), namely that the coefficient of \(L^q\) in \(\bigl(L/(1-e^{-L})\bigr)^{q+1}\) equals \(1\) (put \(q=r-p\)).


